%% 
%% Copyright 2019-2024 Elsevier Ltd
%% 
%% Version 2.4
%% 
%% This file is part of the 'CAS Bundle'.
%% --------------------------------------
%% 
%% It may be distributed under the conditions of the LaTeX Project Public
%% License, either version 1.2 of this license or (at your option) any
%% later version.  The latest version of this license is in
%%    http://www.latex-project.org/lppl.txt
%% and version 1.2 or later is part of all distributions of LaTeX
%% version 1999/12/01 or later.
%% 
%% The list of all files belonging to the 'CAS Bundle' is
%% given in the file `manifest.txt'.
%% 
%% Template article for cas-sc documentclass for 
%% single column output.

%\documentclass[a4paper,fleqn,longmktitle]{cas-sc}
\documentclass[a4paper,fleqn]{cas-sc}

% Author-year citations: \citet{key} -> Author (year),
% \citep{key} -> (Author, year). The `sort' option orders grouped
% parenthetical citations according to the alphabetical bibliography.
\usepackage[authoryear,round,sort]{natbib}
\usepackage{tabularx}
\usepackage[short]{optidef}

\usepackage{graphicx}
\usepackage{float}
\usepackage[section]{placeins}

\renewcommand{\topfraction}{0.9}
\renewcommand{\bottomfraction}{0.8}
\renewcommand{\textfraction}{0.1}
\renewcommand{\floatpagefraction}{0.8}

\setcounter{topnumber}{2}
\setcounter{bottomnumber}{2}
\setcounter{totalnumber}{4}


%%%Author macros
\def\tsc#1{\csdef{#1}{\textsc{\lowercase{#1}}\xspace}}
\tsc{WGM}
\tsc{QE}
\tsc{EP}
\tsc{PMS}
\tsc{BEC}
\tsc{DE}
%%%

\begin{document}
\let\WriteBookmarks\relax
%%%\def\floatpagepagefraction{1}
%%%\def\textpagefraction{.001}


\shorttitle{A B-spline shaping method for low-thrust interplanetary trajectory design}
\shortauthors{J.C. Sanchez}
%\begin{frontmatter}

\title [mode = title]{A B-spline shaping method for low-thrust interplanetary trajectory design}                      
\tnotemark[1,2]

\tnotetext[1]{This document is the results of the research
   project funded by the National Science Foundation.}

\author[1]{Julio C. Sanchez}[type=editor,
                        auid=000,bioid=1,
                        role=Researcher,
                        orcid=0000-0001-0000-0000]
\cormark[1]
\fnmark[1]
\ead{jsanchezm@us.es}
\ead[url]{www.jkkrishnan.in}

\credit{Conceptualization of this study, Methodology, Software}

\affiliation[1]{organization={Universidad de Sevilla, Department of Aerospace Engineering and Fluid Mechanics},
addressline={Camino de los Descubrimientos s/n}, 
city={Sevilla},
            postcode={695013}, 
            state={},
            country={Spain}}



\cortext[cor1]{Corresponding author}
\fntext[fn1]{.}

\begin{abstract}
Shape-based methods remain attractive for preliminary low-thrust trajectory design because they rapidly generate feasible transfers. The shape methods main strength is the dynamics compliance by design. Following the shape method idea, this work presents a cylindrical-coordinate B-spline shaping method for low-thrust interplanetary trajectory design. The proposed method transcribes the spacecraft trajectory in cylindrical coordinates by using a compact B-spline basis and optimizes the accumulated thrust acceleration. The resulting method is an integration-free non-linear program which optimizes the B-spline control coefficients. Initial and endpoint position-velocity, time of flight and number of revolutions are also prescribed in the optimization problem. The results of the B-spline shaping method are compared against the hodographic shaping benchmarks of Gondelach and Noomen for rendezvous transfers from Earth to Mars, near-Earth asteroid 1989ML, comet Tempel 1, and Mercury. On the available \(20\times 15\) grids, the 10-control-point B-spline formulation obtains lower minimum and median \(\Delta V\) than the reproduced higher-order hodographic method in all four benchmark cases, with formal convergence at nearly all grid points. The comparison also shows that B-spline shaping can reduce \(\Delta V\) at the cost of different maximum-thrust behavior, making Pareto analysis essential when thrust acceleration limits are relevant.
\end{abstract}

\begin{graphicalabstract}
\includegraphics{figs/cas-grabs.pdf}
\end{graphicalabstract}

\begin{highlights}
\item Research highlights item 1
\item Research highlights item 2
\item Research highlights item 3
\end{highlights}

\begin{keywords}
quadrupole exciton \sep polariton \sep \WGM \sep \BEC
\end{keywords}


\maketitle


\section{Introduction}
\label{sec:introduction}

The increasing demand for fuel-efficient and high-capability space missions has driven a growing interest in low-thrust propulsion systems, particularly electric thrusters, as alternatives to conventional chemical engines. While chemical propulsion remains indispensable for high-thrust maneuvers — such as launch, planetary orbit insertion, and time-critical trajectory corrections — its fundamental limitation lies in its relatively low specific impulse. Typically, bipropellant systems achieve specific impulses in the range of 300–450 s. By contrast, electric thrusters — including gridded ion engines and Hall-effect thrusters — can reach specific impulses from 1,500 s to over 10,000 s. For a given delta-v budget, this advantage translates directly into a substantially reduced propellant mass requirement through the Tsiolkovsky rocket equation, enabling either leaner spacecraft designs or greater payload capacity. This advantage has been demonstrated extensively in operational missions. The Deep Space 1 spacecraft was the first to validate solar electric propulsion as a primary propulsion system in deep space, successfully completing flybys of asteroid 9969 Braille and comet Borrelly \citet{Rayman2000}. Building on this heritage, the Dawn mission used ion propulsion to become the first spacecraft to orbit two extraterrestrial bodies, achieving orbit at both protoplanet Vesta and dwarf planet Ceres — a mission that would have been infeasible with chemical propulsion within its cost constraints \citet{Brophy2011}. The Japanese Hayabusa and Hayabusa2 missions exploited microwave-discharge and gridded ion engines to accomplish round-trip sample-return missions to near-Earth asteroids Itokawa and Ryugu, respectively \citet{Kuninaka2007, Tsuda2013}. ESA's SMART-1 validated electric primary propulsion for orbit-raising, spiraling from a geostationary transfer orbit to lunar capture over 13 months using only a Hall-effect thruster \citet{Racca2002}. In the commercial sector, Boeing's 702SP satellite platform introduced fully all-electric orbit-raising and stationkeeping using XIPS ion thrusters, eliminating the on-board chemical apogee engine entirely and achieving a dry-mass fraction of approximately 80\% — roughly double that of conventional chemically propelled platforms \citet{Feuerborn2013}. The principal trade-off is thrust magnitude. Electric thrusters produce forces on the order of millinewtons to a few newtons, several orders of magnitude below chemical engines. Accordingly, the operation of electric propulsion is commonly referred to as low-thrust propulsion. Unlike high-thrust maneuvers, which can be adequately modeled as instantaneous velocity changes (impulsive burns), low-thrust trajectories require extended thrusting arcs over many orbital revolutions to achieve the desired orbit change. This renders low-thrust trajectory rendezvous design substantially more complex, motivating the scope of this work.
 
Continuous-thrust trajectory optimization is traditionally formulated as an optimal control problem and solved using Pontryagin’s Minimum Principle \citet{Pontryagin1962}, yielding a nonlinear two-point boundary-value problem in the state and costate variables \citet{Bryson1975}. For bounded-thrust minimum-fuel problems, the optimal throttle follows a bang-bang or bang-off-bang law determined by the switching function (primer vector) \cite{Russell2007, Prussing2010}, making the resulting TPBVP highly sensitive to the initial costate guess and difficult to solve numerically. Due to this reason, low-thrust indirect methods require the addition of heuristic or hybrid techniques to converge \citep{Morante2021}. In \citep{Taheri2016}, an extended logarithmic smoothing technique combined with a continuation strategy is proposed to regularize the bang-bang control structure which improves the convergence of indirect minimum-fuel rendezvous solutions. Similarly, \citet{Tang2018} proposed a hybrid optimization framework combining successive convex programming, pseudospectral methods, and indirect optimization, where adjoint mapping and homotopy continuation are employed to improve costate initialization and robustly converge to fuel-optimal solutions. Likewise, \citet{PerezPalau2018} solved fuel-optimal Earth–Moon transfers via indirect optimal control, relying on a large-scale exploration of the initial costate space to overcome the convergence difficulties of the shooting method. More recently, \citet{Wang2022} developed an indirect framework that combines analytical derivatives, continuation, and switching detection to efficiently solve bang-bang low-thrust problems with interior-point constraints, significantly enlarging the convergence domain while preserving computational efficiency.

Alternatively, direct optimization methods discretize the optimal control problem into a finite nonlinear programming problem \citet{Hargraves1987, Betts1998}, generally providing improved robustness and reducing sensitivity to the initial guess compared with indirect approaches. Although these methods generally sacrifice strict optimality due to the finite-dimensional approximation for the sake of convergence and computational efficiency \citep{Urrios2024}. A popular family of direct methods is the collocation approach \citep{Topputo2014} which approximates the state and control histories by using polynomial functions while enforcing dynamics at a finite set of collocation points. By doing so, the continuous optimal control problem is transformed into a set of algebraic constraints. A popular collocation approach is the pseudospectral method which approximates the state and control histories by using global orthogonal polynomials while enforcing the dynamics at selected collocation points \citep{Garg2010}. An application for low-thrust rendezvous application can be found in \citet{Li2014} which combined shape-based trajectory design with Radau pseudospectral optimization, exploiting the former for rapid preliminary design and the latter for accurate trajectory refinement. Still, direct optimization approaches require an adequate initial guess of the trajectory and control evolution which may be hard for multi-revolution transfers. 

A subclass of direct optimization is the use of shape-based methods in which the spacecraft trajectory or control is parameterized by a family of basis functions. Shape-based methods represent the trajectory with expressive base functions parameterized by small set of design parameters which allow a rapid convergence. The pioneer work in this area stems from \citep{Petropoulos2004} which introduces the exponential sinusoid as a basis function for the orbital radius under tangential thrust. An improvement was made by \citep{NovakVasile2011} by extending the concept to the optimization of the shaping function coefficients. Later, \citep{Gondelach2015} presented the hodographic shaping method which parameterized the velocity in cylindrical coordinates instead of the position evolution. In that publication the tradeoff between a fast low-order, same coefficients as boundary conditions, shape function and an optimized high-order function is detailed. In \citep{Taheri2016} proposes the use of finite Fourier series TBD. More recently, \citep{Wu2022} proposes analytical shaping with a cubic spline function per segment which can be extended for optimality by finding the optimal splines coefficients.

This manuscript proposes a B-spline shaping method for low-thrust trajectory rendezvous. The trajectory is described by cylindrical coordinates which are parameterized with B-splines functions \citep{DeBoor1972}. The cylindrical dynamics has the flatness property \citep{Fliess1995} and the control acceleration can be uniquely determined by the cylindrical coordinates parameterization. The B-spline representation provides a smooth and compact description of the trajectory through a limited number of control points while preserving local support, enabling local trajectory modifications without affecting the entire solution. Consequently, the continuous optimal control problem is transformed into a finite-dimensional nonlinear programming problem in which the B-spline control points constitute the optimization variables. The dynamics is satisfied by design and there are no associated constraints to them. This formulation seeks to balance computational efficiency and trajectory accuracy, making it suitable for the rapid generation of feasible low-thrust rendezvous trajectories.

This work mainly compares with the \citep{Gondelach2015} high-order solutions for transfers to Mars, Mercury, comet Tempel-1 and asteroid 1989ML. Our first main contribution is the demonstration that the optimization of the cylindrical B-splines parameterization produces lower delta-v solutions for the majority of analyzed transfers. The second contribution is that the B-spline support is quite general since only the number of control points and basis polynomial degree have to be adjusted. Differently from \citep{Gondelach2015} there is no need to handcraft the basis functions. Shaped-based B-spline parameterizations have been used in the field of attitude path planning \citep{Louembet2009, Sanchez2020_AA}, model predictive control for cislunar rendezvous \citep{Sanchez2020_AA} and maneuver planning around asteroids \citep{Dambrosio2023}. However, up to the author knowledge, their use has not been explored for low-thrust interplanetary transfers.

The structure of this paper is as follows. Section~\ref{sec:problem} states the low-thrust rendezvous problem in cylindrical coordinates, its polynomial shaping, and the semi-infinite thrust constraint that shaping leaves open. Section~\ref{sec:enforcement} develops the analytical enforcement of that constraint: the Levi-Civita and Kustaanheimo--Stiefel maps that make the heliocentric distance polynomial, the Sundman transformation that lowers the degree of the control numerators, the rational B\'ezier convex-hull condition, and the three shaping families built on it. Section~\ref{sec:implementation} describes the nonlinear program, the seeding procedure, and the post-solve verification. Section~\ref{sec:results} presents numerical results, and Section~\ref{sec:conclusions} concludes the paper.

\section{Shaping-based low-thrust interplanetary rendezvous}
\label{sec:problem}

\subsection{Low-thrust interplanetary rendezvous}
\label{sec:continuous}

This work focuses on the fixed-time rendezvous problem with
low thrust. The transfer time and number of revolutions are fixed, and the goal
is to minimize an objective function which depends on the control. Since the
spacecraft mass is not included as a state variable, the control
$\mathbf{u}$ represents thrust acceleration rather than thrust force. The
trajectory is parameterized in cylindrical coordinates
$\mathbf{q}=[\rho,\,\theta,\,z]^T$, where $\rho$ is the distance of the spacecraft projection
onto the ecliptic plane from the Sun, $\theta$ is the ecliptic longitude, and
$z$ is the displacement normal to the ecliptic. The heliocentric distance is $r=\sqrt{\rho^2+z^2}$. Using the cylindrical velocity $\mathbf{v}=[v_\rho,\,v_\theta,\,v_z]^T$, the
low-thrust rendezvous problem is

\begin{mini}
	{\mathbf{u}(\cdot)}
	{J(\mathbf{u}(\cdot))}
	{\label{prob:main}}{}
	\addConstraint{\dot{\rho}=v_\rho,\qquad\quad\>\>\>
		\dot{v}_\rho=\frac{v_\theta^2}{\rho}-\frac{\mu\rho}{r^3}+u_\rho}
	\addConstraint{\dot{\theta}=\frac{v_\theta}{\rho},\qquad\quad\>\>
		\dot{v}_\theta=-\frac{v_\rho v_\theta}{\rho}+u_\theta}
	\addConstraint{\dot{z}=v_z,\qquad\quad\>\>\>
		\dot{v}_z=-\frac{\mu z}{r^3}+u_z}
	\addConstraint{\mathbf{q}(t_0)=\mathbf{q}_0,\qquad
		\mathbf{v}(t_0)=\mathbf{v}_0}
	\addConstraint{\mathbf{q}(t_f)=\mathbf{q}_f,\qquad
		\mathbf{v}(t_f)=\mathbf{v}_f}
	\addConstraint{\lVert\mathbf{u}(t)\rVert\leq u_{\max},
		\qquad t\in[t_0,t_f].}
\end{mini}

The boundary states are obtained from the ephemerides of the departure and
target bodies at $t_0$ and $t_f$, respectively. The transfer time
$T=t_f-t_0$ and the number of revolutions $N\in\mathbb{N}$ are fixed. The
unwrapped terminal longitude is selected consistently with $N$, i.e., it
differs from the physical terminal longitude by the required multiple of
$2\pi$.

Two objectives are considered. In the minimum-$\Delta V$ problem,
$J=\Delta V=\int_{t_0}^{t_f}\lVert\mathbf{u}(t)\rVert\,\mathrm{d}t$ for a
prescribed acceleration cap $u_{\max}$; solving it for a sequence of caps
traces the $u_{\max}$--$\Delta V$ Pareto front, and when no cap is prescribed
the path constraint is omitted and the peak acceleration is evaluated a
posteriori. In the minimum-peak problem, $u_{\max}$ is itself a decision
variable and $J=u_{\max}$, so the path constraint is active by construction.
In both cases the path constraint must hold at every instant of the transfer,
which is the difficulty addressed in this work.

\subsection{Polynomial shaping as a finite-dimensional problem}
\label{sec:shaping}

Shaping methods parameterize the spacecraft trajectory with a set of basis
functions. Some examples are Fourier series, B\'ezier curves or B-splines. By
parameterizing the trajectory, the dynamics and the boundary conditions can be
satisfied by design. In this work, polynomial shaping (e.g. B\'ezier curves,
B-splines, etc.) is analyzed. Let us parameterize the cylindrical coordinates as
\begin{equation}
	\mathbf{q}(\tau)=\sum^{n_c}_{i=1}\mathbf{c}_i B_i(\tau),\qquad\tau\in[0,1],
	\label{eq:shaping}
\end{equation}
where $B_i$ are polynomial basis functions, $\mathbf{c}_i$ are control points,
and $\tau$ is the curve parameter. Throughout this section $\tau=(t-t_0)/T$ is
the normalized time; Section~\ref{sec:sundman} replaces it by a normalized
Sundman arc, and everything below carries over with the corresponding
change of the endpoint rates. With that parameterization, the derivatives of the
cylindrical coordinates are
\begin{equation}
	\mathbf{q}'(\tau)=\sum^{n_c}_{i=1}\mathbf{c}_i B_i'(\tau),\quad\quad
	\mathbf{q}''(\tau)=\sum^{n_c}_{i=1}\mathbf{c}_i B_i''(\tau),
\end{equation}
where primes denote derivatives with respect to $\tau$, and the physical-time
derivatives follow from the chain rule,
$\dot{\mathbf q}=\mathbf q'/T$ and $\ddot{\mathbf q}=\mathbf q''/T^2$. Then,
the required control acceleration $\mathbf{u}$ can be obtained from the inverse dynamics of
Problem~\eqref{prob:main} as
\begin{equation}
	\mathbf{u}
	=
	\begin{bmatrix}
		u_\rho\\
		u_\theta\\
		u_z
	\end{bmatrix}
	=
	\dfrac{1}{r^3}\begin{bmatrix}
		r^3(\ddot{\rho}-\rho\dot{\theta}^{2})+\mu\rho\\[1mm]
		r^3(2\dot{\rho}\dot{\theta}+\rho\ddot{\theta})\\[1mm]
		r^3\ddot{z}+\mu z
	\end{bmatrix}=\dfrac{1}{r^3}
	\begin{bmatrix}
		P_\rho\\[1mm]
		P_\theta\\[1mm]
		P_z
	\end{bmatrix},
	\label{eq:flatness2}
\end{equation}
so that the control is a rational function of the shaped coordinates and
their first two derivatives. This is the flatness property of the low-thrust cylindrical
dynamics. No integration is required and the control is uniquely determined by
the trajectory shape.

In this work the basis functions $B_i$ are those of a clamped B-spline of
degree $n$. On each knot span the shaped coordinates are polynomials of
degree $n$ and can be converted exactly into Bernstein form using a B-spline
extraction matrix. This local Bernstein representation is subsequently used
to construct polynomial products, integrals, and convex-hull bounds.

Substituting Eq.~\eqref{eq:shaping} into Problem~\eqref{prob:main} gives the
finite-dimensional problem
\begin{mini}
	{\mathbf{C}}
	{J(\mathbf{C})}
	{\label{prob:shaped}}{}
	\addConstraint{\mathbf{q}(0;\mathbf C)=\mathbf{q}_0,\qquad
		\mathbf{q}'(0;\mathbf C)=\mathbf{q}'_0}
	\addConstraint{\mathbf{q}(1;\mathbf C)=\mathbf{q}_f,\qquad
		\mathbf{q}'(1;\mathbf C)=\mathbf{q}'_f}
	\addConstraint{\lVert\mathbf{u}(\tau;\mathbf C)\rVert\leq u_{\max},
		\qquad\tau\in[0,1],}
\end{mini}
where $\mathbf{C}=[\mathbf{c}^T_1,\hdots,\mathbf{c}^T_{n_c}]^T$ stacks all the
control points, which are the decision variables of the problem. The dynamics
are satisfied identically through
Eq.~\eqref{eq:flatness2}, and $\mathbf q'_0$, $\mathbf q'_f$ are the
cylindrical rates obtained from $\mathbf v_0$ and $\mathbf v_f$ and expressed
in the curve parameter. Two ingredients remain to be made algebraic: the
rendezvous conditions, which are linear in the control points and are
eliminated in Section~\ref{sec:rendezvous}, and the thrust constraint, which
is semi-infinite and is the subject of Section~\ref{sec:thrust} and of the
remainder of the paper.

\subsubsection{Rendezvous conditions}
\label{sec:rendezvous}

For a clamped B-spline of degree $n$, only the first basis function is
nonzero at $\tau=0$ and only the first two have a nonzero derivative there;
the same holds at $\tau=1$ for the last two. The endpoint values and rates
therefore determine the two control points adjacent to each end,
\begin{equation}
	\mathbf{c}_1=\mathbf{q}_0,\qquad
	\mathbf{c}_2=\frac{\mathbf{q}'_0-B_1'(0)\,\mathbf{c}_1}{B_2'(0)},\qquad
	\mathbf{c}_{n_c}=\mathbf{q}_f,\qquad
	\mathbf{c}_{n_c-1}
	=\frac{\mathbf{q}'_f-B_{n_c}'(1)\,\mathbf{c}_{n_c}}{B_{n_c-1}'(1)}.
	\label{eq:clamped-elimination}
\end{equation}
The remaining $n_c-4$ control points of each coordinate are the free
variables of Problem~\eqref{prob:shaped}, and the rendezvous conditions are
satisfied exactly by construction. The number of revolutions enters only
through the unwrapped terminal longitude $\theta_f$. When the shaped
quantities are not the cylindrical coordinates themselves but primitives
related to them by a nonlinear map, as in Section~\ref{sec:enforcement}, the
prescribed endpoint values and rates are first transformed to the primitives
and Eq.~\eqref{eq:clamped-elimination} is applied to those.

\subsubsection{Thrust constraint}
\label{sec:thrust}

Let us focus on the control acceleration constraint by expanding the shaped
control acceleration of Eq.~\eqref{eq:flatness2},
\begin{equation}
	\dfrac{1}{r^3}\sqrt{P^2_\rho+P^2_\theta+P^2_z}\leq u_{\max}.
\end{equation}
To simplify the previous expression let us square both sides of the
constraint and multiply the resulting expression by $W^2$,
\begin{equation}
	P^2_\rho+P^2_\theta+P^2_z \leq r^6 u^2_{\max},
	\label{eq:squared-thrust-constraint}
\end{equation}
where the main obstruction to applying polynomial positivity techniques is
that in cylindrical coordinates $r^3=(\rho^2+z^2)^{3/2}$ is not a
polynomial. Note that $r^3$ also appears in $P_\rho$, $P_\theta$ and $P_z$.
Enforcing the inequality at a finite set of nodes is the usual remedy, but it
leaves the constraint unverified between nodes. Section~\ref{sec:enforcement}
removes the obstruction and enforces the bound on the whole trajectory.

\section{Analytical enforcement of the thrust constraint}
\label{sec:enforcement}

\subsection{Levi-Civita map of the meridional plane}
\label{sec:pythagorean}

The obstruction identified in Section~\ref{sec:thrust} is the square root in
$r=\sqrt{\rho^2+z^2}$ which arises when using cylindrical coordinates. For a
fixed $\theta$, the coordinates $(\rho,z)$ describe the spacecraft position
in the meridional plane. To circumvent the obstruction we use the Levi-Civita
map of the meridional plane: writing $\rho+\mathrm{i}z$ as the square of a
complex primitive $w=p+\mathrm{i}q$, the distance becomes
$r=|w^2|=|w|^2$, which is free of square roots. In components, the
Levi-Civita map of the $\rho z$ plane is
\begin{equation}
	\begin{aligned}
		\rho &= p^2-q^2,\\
		z &= 2pq,
	\end{aligned}
	\label{eq:levi-civita-map}
\end{equation}
and the heliocentric distance is
\begin{equation}
	r = \sqrt{\rho^2+z^2}=\sqrt{(p^2+q^2)^2}=p^2+q^2.
	\label{eq:levi-civita-radius}
\end{equation}
Then, we switch from the cylindrical coordinates $(\rho,\>\theta,\> z)$ to the
new set of primitives $(p,\>q,\>\theta)$. The primitives $p$ and $q$ encode
the position within the meridional plane. The map is two-to-one, since
$(p,q)$ and $(-p,-q)$ give the same $(\rho,z)$. The ambiguity is removed by a
single sign choice: the cylindrical description requires $\rho>0$, hence
$p^2>q^2$, so $p$ never vanishes along the trajectory and keeps its sign.
Imposing $p>0$ at the departure point therefore fixes the branch on the
whole arc, and the bound $\rho\geq\rho_{\min}>0$ of the optimization problem
guarantees it during the solution. The longitude $\theta$ is retained as a
third primitive, which keeps the number of revolutions an endpoint value as
in Section~\ref{sec:rendezvous}.

If $p$ and $q$ are polynomials of degree $n$ in the curve parameter, it
follows from Eq.~\eqref{eq:levi-civita-radius} that $r$ is a polynomial of
degree $2n$ and $r^3$ is a polynomial of degree $6n$.

\subsection{Physical-time thrust certificate}
\label{sec:time-certificate}

With the Levi-Civita map of the meridional plane, the thrust constraint
becomes polynomial whenever the primitives are shaped with polynomial basis
functions. The time derivatives of $\rho$ and $z$ are
\begin{equation}
	\begin{alignedat}{2}
		\dot \rho
		&=2(p\dot p-q\dot q),
		&\qquad
		\dot z
		&=2(\dot p\,q+p\dot q),
		\\
		\ddot \rho
		&=2\left(\dot p^{2}+p\ddot p
		-\dot q^{2}-q\ddot q\right),
		&\qquad
		\ddot z
		&=2\left(\ddot p\,q+2\dot p\dot q+p\ddot q\right).
	\end{alignedat}
	\label{eq:levi-civita-time-derivatives}
\end{equation}
Let us define the kinematic acceleration terms
\begin{equation}
	\begin{aligned}
		\mathcal A_\rho
		&=\ddot{\rho}-\rho\dot{\theta}^2=2\left(\dot p^{2}+p\ddot p
		-\dot q^{2}-q\ddot q\right)-(p^2-q^2)\dot{\theta}^2,\\
		\mathcal A_\theta
		&=2\dot{\rho}\dot{\theta}+\rho\ddot{\theta}=4(p\dot p-q\dot q)\dot{\theta}+(p^2-q^2)\ddot{\theta},\\
		\mathcal A_z
		&=\ddot{z}=2\left(\ddot p\,q+2\dot p\dot q+p\ddot q\right),
	\end{aligned}
	\label{eq:time-kinematic-terms}
\end{equation}
and the cubic heliocentric distance $r^3$ as
\begin{equation}
	r^3 = (p^2+q^2)^3=p^6+3p^4q^2+3p^2q^4+q^6.
	\label{eq:levi-civita-weight}
\end{equation}
Then, the thrust constraint of Eq.~\eqref{eq:squared-thrust-constraint} can
be written in the meridional Levi-Civita primitives and the heliocentric
longitude as
\begin{equation}
	\underbrace{
		\left(r^3\mathcal A_\rho+\mu (p^2-q^2)\right)^2
	}_{P_\rho^{2}}
	+
	\underbrace{
		\left(r^3\mathcal A_\theta\right)^2
	}_{P_\theta^{2}}
	+
	\underbrace{
		\left(r^3\mathcal A_z+2\mu pq\right)^2
	}_{P_z^{2}}
	\leq
	r^6\,u_{\max}^2.
	\label{eq:pythagorean-time-certificate}
\end{equation}
Time derivatives are taken with respect to the normalized time $\tau$, which
only introduces powers of the fixed transfer time $T$ and does not change the
polynomial degrees. If $p$, $q$, and $\theta$ are degree-$n$ polynomials on
each span, the degrees of the terms in
Eq.~\eqref{eq:pythagorean-time-certificate} are given in
Table~\ref{tab:thrust-certificate-degrees}.

\begin{table}[htbp]
	\centering
	\caption{Polynomial degree of the thrust constraint for degree-$n$
	$(p,q,\theta)$ parameterized with time.}
	\label{tab:thrust-certificate-degrees}
	\begin{tabular}{lc}
		\toprule
		\textbf{Polynomial term} & \textbf{Degree} \\
		\midrule
		$P_\rho^{2}$   & $20n-4$ \\
		$P_\theta^{2}$ & $18n-4$ \\
		$P_z^{2}$      & $16n-4$ \\
		$r^6$      & $12n$   \\
		\midrule
		$r^6 u_{\max}^2-P_\rho^{2}-P_\theta^{2}-P_z^{2}$
		& $20n-4$ \\
		\bottomrule
	\end{tabular}
\end{table}

Requiring every Bernstein coefficient of the degree-$(20n-4)$ polynomial
$r^6u_{\max}^2-P_\rho^{2}-P_\theta^{2}-P_z^{2}$ to be nonnegative on each span
is a sufficient condition for the thrust constraint and constitutes the
physical-time baseline of this work. Its degree, however, makes the condition
expensive and ill-conditioned already for moderate $n$, which motivates the
two developments that follow: a change of independent variable that lowers
the degree of the numerators below that of $r^3$
(Section~\ref{sec:sundman}), and a convex-hull condition that never forms the
squared polynomial (Section~\ref{sec:hull}).

\subsection{Sundman transformation}
\label{sec:sundman}

The Sundman transformation reduces the polynomial degree that arises when the
primitives are shaped in physical time $t$. It replaces $t$ with the parameter
$s$ through
\begin{equation}
	\mathrm{d}t=r\,\mathrm{d}s,
	\label{eq:sundman-transformation}
\end{equation}
where $r>0$. Denoting physical-time derivatives with dots and derivatives
with respect to the Sundman parameter $s$ with primes, the chain rule gives
\begin{equation}
	\dot f=\frac{f'}{r},
	\qquad
	\ddot f
	=\frac{1}{r}\left(\frac{f'}{r}\right)'
	=\frac{rf''-r'f'}{r^3}.
	\label{eq:sundman-chain-rule}
\end{equation}
The second derivative carries the same denominator $r^3$ as the gravity
term, so no further multiplication by $r^3$ is needed to clear the control.
In cylindrical components, substituting Eq.~\eqref{eq:sundman-chain-rule}
into Eq.~\eqref{eq:flatness2} gives
\begin{equation}
	\mathbf{u}=\dfrac{1}{r^3}\begin{bmatrix}
		r\rho''-r'\rho'-r\rho\,(\theta')^2+\mu\rho\\
		\rho\,(r\theta''-r'\theta')+2r\rho'\theta'\\
		rz''-r'z'+\mu z
	\end{bmatrix},
	\label{eq:sundman-control}
\end{equation}
Each entry equals $r^3$ times the physical
kinematic acceleration of Eq.~\eqref{eq:time-kinematic-terms} plus the
gravity numerator, but the factor $r^3$ is now supplied by the chain rule
instead of being multiplied in. Accordingly, in the Sundman parameter the
kinematic terms are $r^3$ times the physical ones
\begin{equation*}
\mathcal{B}_\rho=r^3\mathcal{A}_\rho,\quad\mathcal{B}_\theta=r^3\mathcal{A}_\theta,\quad \mathcal{B}_z=r^3\mathcal{A}_z,
\end{equation*}
which are also epxressed as
\begin{equation}
	\begin{aligned}
		\mathcal B_\rho
		&=r\rho''-r'\rho'-r\rho\,(\theta')^2
		=2r\left(pp''-qq''\right)
		-\rho\left[2\left(p'^2+q'^2\right)+r\,(\theta')^2\right],\\
		\mathcal B_\theta
		&=\rho\,(r\theta''-r'\theta')+2r\rho'\theta'
		= r\rho\,\theta''
		+2\left[pp'\left(p^2+3q^2\right)-qq'\left(3p^2+q^2\right)\right]\theta',\\
		\mathcal B_z
		&=rz''-r'z'
		= 2r\left(qp''+pq''\right)-2z\left(p'^2+q'^2\right),
	\end{aligned}
	\label{eq:sundman-kinematic-pq}
\end{equation}
with $r=p^2+q^2$, $\rho=p^2-q^2$ and $z=2pq$. The thrust constraint of
Eq.~\eqref{eq:squared-thrust-constraint} becomes
\begin{equation}
	\underbrace{
		\left(\mathcal B_\rho+\mu (p^2-q^2)\right)^2
	}_{P_\rho^{2}}
	+
	\underbrace{
		\mathcal B_\theta^{2}
	}_{P_\theta^{2}}
	+
	\underbrace{
		\left(\mathcal B_z+2\mu pq\right)^2
	}_{P_z^{2}}
	\leq
	r^6\,u_{\max}^2.
	\label{eq:sundman-certificate}
\end{equation}
Each numerator contains at most one factor of $r$ together with two
derivatives, so $P_\rho$, $P_\theta$ and $P_z$ have degrees $6n-2$, $5n-2$
and $4n-2$, all below $\deg W=6n$. The squared constraint of
Eq.~\eqref{eq:sundman-certificate} has degree $12n$ instead of the $20n-4$
of Eq.~\eqref{eq:pythagorean-time-certificate}.

\begin{table}[htbp]
	\centering
	\caption{Polynomial degree of the thrust constraint for degree-$n$
		$(p,q,\theta)$ parameterized with the Sundman parameter.}
	\label{tab:sundman-certificate-degrees}
	\begin{tabular}{lc}
		\toprule
		\textbf{Polynomial term} & \textbf{Degree} \\
		\midrule
		$P_\rho^{2}$   & $12n-4$ \\
		$P_\theta^{2}$ & $10n-4$ \\
		$P_z^{2}$      & $8n-4$ \\
		$r^6$      & $12n$   \\
		\midrule
		$r^6 u_{\max}^2-P_\rho^{2}-P_\theta^{2}-P_z^{2}$
		& $12n$ \\
		\bottomrule
	\end{tabular}
\end{table}

The change of independent variable alters the curve parameter of
Section~\ref{sec:shaping}. The primitives are now shaped in the normalized
Sundman arc $\sigma=(s-s_0)/S\in[0,1]$, where the total arc
$S=s_f-s_0$ is unknown a priori and becomes an additional decision variable;
derivatives with respect to $s$ are those with respect to $\sigma$ divided by
$S$. The rendezvous conditions of Section~\ref{sec:rendezvous} are imposed
on the primitives in $\sigma$. To preserve the prescribed physical endpoint
velocities, the primitive endpoint derivatives must satisfy
\begin{equation}
	\frac{\mathrm{d}f}{\mathrm{d}\sigma}=S\,r\,\dot f
	\label{eq:sundman-endpoint-rates}
\end{equation}
for every primitive $f$, where $r$ is known at the endpoints, so the
eliminated control points of Eq.~\eqref{eq:clamped-elimination} are affine
in $S$. The prescribed time of flight is enforced through
\begin{equation}
	T=t_f-t_0
	=\int_{s_0}^{s_f}r\,\mathrm{d}s
	=S\int_{0}^{1}r(\sigma)\,\mathrm{d}\sigma,
	\label{eq:normalized-sundman-time-of-flight}
\end{equation}
whose integrand is polynomial on each span, so the constraint is polynomial
in the control points and linear in $S$.

\subsection{Rational B\'ezier convex-hull enforcement}
\label{sec:hull}

On each span, with local parameter $\sigma\in[0,1]$, any polynomial of degree
at most $m$ can be written in the Bernstein basis
$B^m_i(\sigma)=\binom{m}{i}\sigma^i(1-\sigma)^{m-i}$. Let $m$ be the largest
of the degrees of $P_\rho$, $P_\theta$, $P_z$ and $r^3$, given in
Table~\ref{tab:hull-degrees}, and let the remaining polynomials be
degree-elevated to $m$, which is exact and linear in their coefficients.
Then $\mathbf{P}=[P_{\rho},\>P_{\theta},\>P_z]^T$ and $r^3$ share the same
basis,
\begin{equation}
	\mathbf{P}(\sigma)=\sum^{m}_{i=0}\mathbf{P}_i B^m_i(\sigma),\qquad
	r^3(\sigma)=\sum^{m}_{i=0}W_i B^m_i(\sigma),
\end{equation}
and the control is the ratio
\begin{equation}
	\mathbf{u}(\sigma)=\dfrac{\mathbf{P}(\sigma)}{r^3(\sigma)}
	=\dfrac{\sum\limits^{m}_{i=0}\mathbf{P}_i B^m_i(\sigma)}
	{\sum\limits^{m}_{j=0}W_j B^m_j(\sigma)}.
\end{equation}
Defining the control points $\mathbf{c}_i=\mathbf{P}_i/W_i$, so that
$\mathbf{P}_i=W_i\mathbf{c}_i$, and splitting the fraction term by term,
\begin{equation}
	\mathbf{u}(\sigma)
	=\sum^{m}_{i=0}\lambda_i(\sigma)\,\mathbf{c}_i,
	\qquad
	\lambda_i(\sigma)=\dfrac{W_i B^m_i(\sigma)}
	{\sum\limits^{m}_{j=0}W_j B^m_j(\sigma)},
	\label{eq:rational-bezier}
\end{equation}
which is a rational B\'ezier curve with control points $\mathbf{c}_i$ and
weights $W_i$.

By construction $\sum_i\lambda_i(\sigma)=1$. Since $B^m_i(\sigma)\geq 0$ on
$\sigma\in[0,1]$ and $\sum_j B^m_j(\sigma)=1$, positive weights $W_i>0$ give
$\sum_j W_jB^m_j(\sigma)\geq\min_j W_j>0$ and $\lambda_i(\sigma)\geq 0$.
The control is then a convex combination of the control points and lies
inside the convex hull of $\{\mathbf{c}_0,\ldots,\mathbf{c}_m\}$ for every
$\sigma\in[0,1]$.

Suppose now that $\lVert\mathbf{c}_i\rVert\leq u_{\max}$ for
$i=0,\ldots,m$. By the triangle inequality,
\begin{equation*}
	\lVert\mathbf{u}(\sigma)\rVert
	=\Big\lVert\sum_i\lambda_i(\sigma)\mathbf{c}_i\Big\rVert
	\leq\sum_i\lambda_i(\sigma)\lVert\mathbf{c}_i\rVert
	\leq u_{\max},
	\qquad \sigma\in[0,1].
\end{equation*}
Since $\mathbf{c}_i=\mathbf{P}_i/W_i$ and $W_i>0$, the thrust constraint on
the whole span is therefore guaranteed by the $m+1$ conditions
\begin{equation}
	P^2_{\rho,i}+P^2_{\theta,i}+P^2_{z,i}\leq u^2_{\max}W^2_i,
	\qquad i=0,\ldots,m,
	\label{eq:hull}
\end{equation}
written in squared form to avoid both the division by $W_i$ and the
nonsmooth norm. The positivity $W_i>0$ is inherited from the lower bound
imposed on the Bernstein coefficients of $r$, since products and degree
elevation of positive Bernstein coefficients remain positive.

The condition is sufficient but not necessary, because the interior control
points do not lie on the curve; subdividing the span brings them closer to
it and reduces the conservatism.

\begin{table}[htbp]
	\centering
	\caption{Polynomial degrees entering the convex-hull condition for
		degree-$n$ $(p,q,\theta)$, parameterized with time and with the
		Sundman parameter.}
	\label{tab:hull-degrees}
	\begin{tabular}{lcc}
		\toprule
		\textbf{Polynomial term} & \textbf{Time} & \textbf{Sundman} \\
		\midrule
		$P_\rho$   & $10n-2$ & $6n-2$ \\
		$P_\theta$ & $9n-2$  & $5n-2$ \\
		$P_z$      & $8n-2$  & $4n-2$ \\
		$r^3$      & $6n$    & $6n$   \\
		\midrule
		Common degree $m$             & $10n-2$ & $6n$   \\
		Inequalities per span, $m+1$  & $10n-1$ & $6n+1$ \\
		\bottomrule
	\end{tabular}
\end{table}



\section{Numerical implementation}
\label{sec:implementation}

\subsection{Nonlinear program}

On each family the decision variables are the interior control points of
every primitive, the Sundman arc length $S$ that maps the spline parameter
$\sigma\in[0,1]$ to $s=S\sigma$, and, in the minimum-peak problem, the squared
bound $u_{\max}^2$. The control points adjacent to each clamped end are
eliminated through Eq.~\eqref{eq:clamped-elimination}. Since the primitive
Sundman rates are $f'=S\,r\dot f$ and $r$ is known at the endpoints, these
control points are affine in $S$, so the elimination remains analytical when
$S$ is free. The constraints are the hull inequalities of
Eq.~\eqref{eq:rational-hull-constraints} on every subdivided block, the
coefficient bounds $r_j\geq r_{\min}$ and, for the cylindrical families,
$\rho_i\geq\rho_{\min}$, the time-of-flight equality of
Eq.~\eqref{eq:normalized-sundman-time-of-flight}, which is bilinear in $S$ and
the coefficients of $r$, and, for the KS family with the revolution lock, the
linear half-plane inequalities of Eq.~\eqref{eq:ks-revolution-lock}. The
hull inequalities are made dimensionless with a reference squared
acceleration and the sixth power of the larger endpoint radius.

In the minimum-peak problem the objective is $u_{\max}^2$, normalized by the
hull bound of the seed trajectory. In the minimum-$\Delta V$ problem the
objective is $S\int_0^1\lVert\mathbf u\rVert\,r\,\mathrm{d}\sigma$, evaluated
by Gauss--Legendre quadrature on every span; the nonsmooth norm is replaced at
each node by an auxiliary bound $\eta_k\geq\lVert\mathbf u_k\rVert$, so that
the program stays smooth. A prescribed cap enters through the hull
inequalities with fixed $u_{\max}$; without a cap no thrust inequality is
imposed and the hull bound of the solution is reported a posteriori. A
quadratic objective $S\int_0^1\lVert\mathbf u\rVert^2 r\,\mathrm{d}\sigma$
with a fixed cap is also available. The program is assembled symbolically
with CasADi and solved with IPOPT, with exact first and second derivatives
obtained by algorithmic differentiation.

\subsection{Seeding}

All families are initialized from the same minimum-$\Delta V$ cylindrical
B-spline solution in physical time. The seed trajectory is sampled, the
Sundman parameter is obtained from $s(t)=\int\mathrm{d}t/r$, and the sampled
path is lifted to the primitives of the family: by the half-angle relations
for the cylindrical families and by the continuous fibre lift for the KS
family. The primitive splines are then fitted in $\sigma$ by least squares
with the endpoint values and rates matched exactly, and $S$ is taken from the
time-of-flight identity. The peak variable is initialized at the hull bound
of the fitted seed times a small factor, and the lifted seed also supplies
the half-plane directions of the revolution lock. Because the seed is a
minimum-$\Delta V$ shape, its hull bound is in general far above the
minimum-peak optimum, and the solver is started from a feasible point whose
thrust constraint is inactive.

\subsection{Post-solve verification}

Every returned solution is checked independently of the solver. The
Bernstein coefficients of $\mathbf P$ and $W$ are recomputed from the control
points, and the solution is accepted only if all weights are strictly
positive, the hull residuals are nonnegative up to a small tolerance, the
control norm sampled densely in $\sigma$ does not exceed the certified bound,
the transfer time and the endpoint positions and velocities match the
prescribed values, and the identity $r^2=\rho^2+z^2$ holds to roundoff. The
reported peak acceleration is the hull bound of the returned control points,
not the value of the IPOPT variable. The hull condition itself is exact, but
the coefficients are evaluated in floating point and the active inequalities
are satisfied only to the solver tolerance, so the achieved bound must be
interpreted together with the actual constraint residuals; the checks above
are consistency tests, not a formal interval certificate.

\section{Numerical results}
\label{sec:results}

% TODO: lead paragraph: transfers (Mars, Mercury, ...), units, hardware.

\subsection{Hull versus node-wise enforcement}

% TODO: minimum-peak solve of one transfer with the hull condition and with
% node-wise limiting at 15 and 25 nodes per span; sampled peak above the
% nodal limit between nodes; figure hull_vs_grid.

\subsection{Comparison of shaping families}

% TODO: Levi-Civita, scaled Levi-Civita, KS (with and without the revolution lock)
% on the same seeds; hull degree, rows per span, iterations, wall time,
% failures; minimum-peak and capped minimum-Delta V.

\subsection{Rendezvous transfers}

% TODO: minimum-peak and Pareto results against the cylindrical B-spline and
% hodographic benchmarks.

\section{Conclusions}
\label{sec:conclusions}

% TODO.

\section*{CRediT authorship contribution statement}

Julio Author: Conceptualization, Methodology, Software, Validation, Writing - original draft. Second Author: Supervision, Writing - review and editing.

\section*{Declaration of competing interest}

The authors declare that they have no known competing financial interests or personal relationships that could have appeared to influence the work reported in this paper.

\section*{Data and code availability}

The comparison scripts and generated CSV files are contained in the FASTTRANSFER repository used to prepare this manuscript. The principal scripts are \path{experiments/compare_gondelach_fig2_fig3_bspline10.py} and \path{experiments/plot_pareto_from_csv.py}.

\section*{Acknowledgements}

Funding, institutional support, and computational-resource acknowledgements will be added before submission.

\bibliographystyle{cas-model2-names}
\bibliography{cas-refs.bib}



\end{document}
